\documentclass{article}
\usepackage[preprint]{neurips_2024}

%% ── Package imports ──────────────────────────────────────────────────────────
\usepackage{amsmath,amssymb,amsthm}
\usepackage{mathtools}
\usepackage{graphicx}
\usepackage{booktabs}
\usepackage{multirow}
\usepackage{array}
\usepackage{xcolor}
\usepackage{algorithm}
\usepackage{algpseudocode}
\usepackage{cleveref}

%% ── Custom maths ─────────────────────────────────────────────────────────────
\newcommand{\Msem}{M_{\mathrm{sem}}}
\newcommand{\Mepi}{M_{\mathrm{epi}}}
\newcommand{\ks}{\mathbf{k}_s}
\newcommand{\ke}{\mathbf{k}_e}
\newcommand{\ksnorm}{\hat{\mathbf{k}}_s}
\newcommand{\kenorm}{\hat{\mathbf{k}}_e}
\newcommand{\vps}{v_{p,s}}
\newcommand{\vpe}{v_{p,e}}
\newcommand{\bx}{\mathbf{x}}
\newcommand{\br}{\mathbf{r}}
\newcommand{\bq}{\mathbf{q}}

%% ── Placeholders ─────────────────────────────────────────────────────────────
\providecommand{\todo}[1]{\textcolor{red}{\textbf{[TODO: #1]}}}

% ─────────────────────────────────────────────────────────────────────────────

\title{DREx: Episodic/Semantic Associative Memory with\\
       Length-Adaptive EMA for Long-Context Transformers}

\author{%
  Wesley Scholl \\
  Independent Researcher \\
  \texttt{wesleyscholl@github} \\
}

\begin{document}

\maketitle

\begin{abstract}
We present \textbf{DREx}, a transformer architecture augmented with a four-tier memory
hierarchy whose L4 layer is a novel episodic/semantic split delta-rule associative memory
module.  The L4 module maintains two independent $\tfrac{H}{2} \times \tfrac{H}{2}$
associative matrices — one for uniformly-weighted semantic associations, one for
recency-weighted episodic bindings — updated via an outer-product delta rule under an
OR relative-norm write gate.  We identify and resolve a previously undocumented failure
mode of EMA-based associative memory at short sequence lengths (the \emph{EMA bootstrap
problem}): a fixed decay coefficient $\alpha$ calibrated for long sequences produces
a memory time-constant $\tau = 1/(1-\alpha)$ that is a large fraction of the sequence
length $L$, causing the write gate to over-fire at $\mathrm{wr} \approx 0.97$.  We
derive a length-adaptive formula $\alpha(L) = 0.95^{96/L}$ that holds $\tau/L \approx
0.21$ constant across $L=32$--$128$.  We further characterize OR-gate write-rate
inflation when two partially-correlated branches each fire with probability $p$, and
confirm a recalibrated threshold $\mathrm{thresh}^* = 0.70$ that delivers validated
write rates $\mathrm{wr}(L{=}32) = 0.581$ and $\mathrm{wr}(L{=}96) = 0.421$.
All architectural decisions are grounded in 247+ controlled ablation experiments over
15 research phases with $\geq\!2/3$ seed confirmation per finding.
Ablations confirm a null retrieval gate is essential ($+0.30$ perplexity without it).
End-to-end evaluation on TinyStories language modelling and long-context passkey/BABILong
benchmarks is in progress; preliminary 2\,000-step probes give validation perplexity
4.27 (memory) vs.\ 4.21 (baseline).
\end{abstract}

% ─────────────────────────────────────────────────────────────────────────────
\section{Introduction}
\label{sec:intro}
% ─────────────────────────────────────────────────────────────────────────────

The capacity of standard transformer attention to retain information beyond its context
window is a central limitation for long-document language modelling, multi-turn dialogue,
and structured reasoning.  Several lines of work address this by replacing or augmenting
soft attention with a compressed, updateable memory store: cross-segment delta-rule
matrices \citep{munkhdalai2024infini}, selective-state recurrences
\citep{gu2024mamba,peng2023rwkv,sun2023retnet}, and test-time weight adaptation
\citep{sun2024titans}.

Associative matrices (Hopfield-style stores \citep{ramsauer2021hopfield} updated by a
delta rule) are attractive because they support $O(1)$ parameter operations per token
and can be inserted as a residual branch inside each transformer layer without altering
the attention stack.  Their write dynamics, however, are sensitive to the relationship
between the EMA decay coefficient $\alpha$, the sequence length $L$, and the stimulus
novelty distribution.  This paper documents, formally characterises, and resolves two
failure modes that we encountered during development:

\begin{enumerate}
\item \textbf{EMA bootstrap failure at short $L$} (\S\ref{sec:ema}).  A fixed
      $\alpha = 0.95$ calibrated for $L=96$ makes the memory time-constant
      $\tau = 20$ steps, i.e.\ $\tau/L = 0.21$, which is well-behaved.  At $L=32$,
      however, the same $\alpha$ gives $\tau/L = 0.625$: the memory never fully decays
      within a single context, every token appears novel to the gate, and
      $\mathrm{wr} \to 0.97$.  We derive a closed-form fix.

\item \textbf{OR-gate write-rate inflation} (\S\ref{sec:gate}).  Splitting memory into
      two branches (semantic + episodic) and OR-ing their write decisions inflates the
      true write rate above the per-branch calibration point.  If each branch fires with
      probability $p$ independently, the OR fires at $1-(1-p)^2$; with correlation the
      inflation is less but still substantial.  We characterise the effect and confirm a
      per-architecture threshold $\mathrm{thresh}^* = 0.70$.
\end{enumerate}

Beyond these engineering contributions, the paper describes the complete validated
architecture of the L4 MemoryModule: the episodic/semantic split, recency weighting,
the null retrieval gate, and the output LayerNorm.  We report every component's evidence
level (high / medium) together with the full list of design choices that were
experimentally ruled out.

\paragraph{Contributions.}
\begin{itemize}
  \item A four-tier memory transformer (\S\ref{sec:arch}) with a novel L4
        episodic/semantic associative layer integrated as a per-layer residual branch.
  \item Identification and analytic treatment of the \emph{EMA bootstrap problem}, with
        the length-adaptive fix $\alpha(L) = 0.95^{96/L}$
        (\S\ref{sec:ema}).
  \item Characterisation of OR-gate write-rate inflation and an experimentally confirmed
        threshold $\mathrm{thresh}^*=0.70$ (\S\ref{sec:gate}).
  \item A comprehensive ablation record  ($247+$ experiments across 15 phases) that
        catalogues both successful components and ruled-out alternatives
        (\S\ref{sec:ablations}).
  \item An open-source implementation with a 199-test suite at 100\% branch coverage
        (\S\ref{sec:impl}).
\end{itemize}

% ─────────────────────────────────────────────────────────────────────────────
\section{Related Work}
\label{sec:related}
% ─────────────────────────────────────────────────────────────────────────────

\paragraph{Associative memory in transformers.}
Modern Hopfield networks \citep{ramsauer2021hopfield} reinterpret attention as a
continuous Hopfield energy minimisation.  Infini-Attention \citep{munkhdalai2024infini}
adds a compressive delta-rule associative segment memory (our L2) alongside local
attention, enabling infinite-length processing.  DREx extends this approach with a
separate per-segment episodic/semantic split module (our L4) that is written
selectively via a relative-norm gate and read as a residual at the query position only.

\paragraph{Linear recurrences and state-space models.}
Mamba \citep{gu2024mamba}, RWKV \citep{peng2023rwkv}, and RetNet
\citep{sun2023retnet} achieve $O(L)$ memory via exponentially decaying state.  Our L2
weight matrix is analogous to the hidden state of these models; the novelty in DREx is
the \emph{additional} per-segment selective write layer (L4) with a structurally
different update rule.

\paragraph{Test-time weight adaptation.}
Titans \citep{sun2024titans} memorises information by gradient-updating a small MLP's
weights at test time.  Our L3 layer is inspired by this idea (Rust-backed weight
snapshots with async prefetch); L4 is complementary — a faster, in-activation store.

\paragraph{Long-context benchmarks.}
We evaluate on passkey retrieval \citep{munkhdalai2024infini} and BABILong
\citep{schalk2024babilong}, both standard benchmarks for long-context recall.  Training
uses TinyStories \citep{eldan2023tinystories} as the language modelling corpus.

% ─────────────────────────────────────────────────────────────────────────────
\section{Architecture}
\label{sec:arch}
% ─────────────────────────────────────────────────────────────────────────────

\subsection{Memory Hierarchy}
\label{sec:hierarchy}

DREx organises memory across four tiers with increasing latency and capacity
(\Cref{tab:hierarchy}).  All four tiers are active simultaneously; each addresses a
different temporal scope.

\begin{table}[h]
  \centering
  \small
  \caption{DREx four-tier memory hierarchy.  L4 is the primary contribution of this paper.}
  \label{tab:hierarchy}
  \begin{tabular}{clll}
    \toprule
    \textbf{Tier} & \textbf{Mechanism} & \textbf{Scope} & \textbf{Location} \\
    \midrule
    L1 & Sliding-window causal attention & In-context (short range) & Activations \\
    L2 & Infini-Attention delta-rule matrix \citep{munkhdalai2024infini}
       & Cross-segment & Activations \\
    L3 & Titans-style MLP weight snapshot \citep{sun2024titans}
       & Disk (long range, async) & Rust SnapshotStore \\
    L4 & Episodic/semantic split delta-rule (\S\ref{sec:l4})
       & Per-segment associative & Activations \\
    \bottomrule
  \end{tabular}
\end{table}

L1 provides standard in-context attention over a sliding window.  L2 compresses
cross-segment information using Infini-Attention's delta-rule update with the
$\phi(\bx) = \mathrm{ELU}(\bx)+1$ feature map, accumulating a normalisation vector
$\mathbf{z}$ for stable reads.  L3 stores slow weight snapshots asynchronously to an
on-disk Rust-backed store with a sketch-based $k$-NN prefetch index.  L4 is the focus
of this paper.

\subsection{L4 MemoryModule}
\label{sec:l4}

The L4 MemoryModule is inserted as an optional residual branch in each transformer
layer.  Its contribution is added at the query (last-token) position only:
\begin{equation}
  \bx_{L-1} \leftarrow \bx_{L-1} + \mathrm{MemoryModule}\!\left(\mathrm{LN}(\bx)\right),
\end{equation}
where $\mathrm{LN}$ denotes pre-LayerNorm \citep{xiong2020layer} applied to the full
context tensor $\bx \in \mathbb{R}^{B \times L \times H}$.

\paragraph{Associative matrices.}
We maintain two independent $\tfrac{H}{2} \times \tfrac{H}{2}$ matrices per forward
call, initialised to zero:
\begin{align}
  \Msem &\in \mathbb{R}^{B \times H/2 \times H/2}, \\
  \Mepi &\in \mathbb{R}^{B \times H/2 \times H/2}.
\end{align}
The fixed 50/50 split was confirmed across 9 seeds (exp\_38\_1); a learned router
reduced accuracy by 10--24\% (\S\ref{sec:ablations}).

\paragraph{Key projections.}
Two separate linear projections (no bias) map each context position $\bx_t \in
\mathbb{R}^{B \times H}$ to sub-space keys:
\begin{equation}
  \ks = W_s \bx_t \in \mathbb{R}^{B \times H/2},
  \qquad
  \ke = W_e \bx_t \in \mathbb{R}^{B \times H/2}.
\end{equation}

\paragraph{Delta-rule write.}
Unit-normalised keys are used throughout to decouple magnitude from direction:
\begin{equation}
  \ksnorm = \frac{\ks}{\|\ks\|},
  \qquad
  \kenorm = \frac{\ke}{\|\ke\|}.
\end{equation}
The current memory predictions for position $t$ are:
\begin{equation}
  \vps = \Msem \ksnorm,
  \qquad
  \vpe = \Mepi \kenorm.
\end{equation}
The outer-product delta rule computes prediction errors and updates:
\begin{align}
  \Delta\Msem &= (\ks - \vps) \otimes \ksnorm, \label{eq:delta-sem} \\
  \Delta\Mepi &= (\ke - \vpe) \otimes \kenorm. \label{eq:delta-epi}
\end{align}

\paragraph{Recency weighting.}
Episodic memory applies position-proportional recency weighting: later observations
are written with higher weight than earlier ones within the same context:
\begin{equation}
  w_t = \frac{t+1}{L} \in \left(0, 1\right].
\end{equation}

\paragraph{Write gate.}
A relative-norm criterion determines whether position $t$ carries novel information
(see \S\ref{sec:gate} for derivation):
\begin{equation}
  g_t =
  \mathbf{1}\!\left[\|\ks - \vps\| \geq \tau^* \|\ks\|
  \;\vee\;
  \|\ke - \vpe\| \geq \tau^* \|\ke\|\right],
  \qquad \tau^* = 0.70.
  \label{eq:gate}
\end{equation}
OR rather than AND was confirmed superior (AND degrades recall; exp\_47\_2).

\paragraph{EMA update.}
Matrices are updated via an exponential moving average with length-adaptive coefficient
$\alpha(L) = 0.95^{96/L}$ (derivation in \S\ref{sec:ema}):
\begin{align}
  \Msem &\mathrel{+}= (1 - \alpha(L))\cdot g_t \cdot \Delta\Msem,
  \label{eq:ema-sem} \\
  \Mepi &\mathrel{+}= (1 - \alpha(L))\cdot w_t \cdot g_t \cdot \Delta\Mepi.
  \label{eq:ema-epi}
\end{align}

\paragraph{Read and output.}
After processing positions $t = 0, \ldots, L-2$, the query position $L-1$ reads from
both matrices:
\begin{equation}
  \br = \left[\Msem \hat{\mathbf{q}}_s \,;\, \Mepi \hat{\mathbf{q}}_e\right]
  \in \mathbb{R}^{B \times H},
  \label{eq:read}
\end{equation}
where $\hat{\mathbf{q}}_s, \hat{\mathbf{q}}_e$ are unit-normalised projections of
$\bx_{L-1}$.  Hard concatenation (no learned read gate) was confirmed optimal (exp\_38\_3;
a learned gate reduced accuracy by 10\%).

A learned null retrieval gate suppresses reads when both matrices are empty:
\begin{equation}
  \tilde{\br} = \sigma\!\left(w^\top \bx_{L-1} + b\right) \cdot \br.
  \label{eq:null-gate}
\end{equation}
Removing this gate costs $+0.30$ validation perplexity (\S\ref{sec:ablations}).

The final output is projected and normalised:
\begin{equation}
  \mathrm{out} = \mathrm{LN}\!\left(W_o \,\tilde{\br}\right) \in \mathbb{R}^{B \times H}.
  \label{eq:output}
\end{equation}
The output LayerNorm bounds the residual contribution when write-path gradients are
absent (detached write optimisation; \S\ref{sec:impl}).

The complete forward procedure is summarised in \Cref{alg:memory}.

\begin{algorithm}[t]
  \caption{L4 MemoryModule forward pass}
  \label{alg:memory}
  \begin{algorithmic}[1]
    \Require $\bx \in \mathbb{R}^{B \times L \times H}$, thresh $\tau^*$, decay $\alpha$
    \Ensure output $\in \mathbb{R}^{B \times H}$ (residual for position $L-1$)
    \State $\Msem, \Mepi \leftarrow \mathbf{0}$  \Comment{zero-initialise per call}
    \For{$t = 0$ to $L-2$}
      \State $\ks \leftarrow W_s\bx_t$; $\ke \leftarrow W_e\bx_t$
      \State $\ksnorm \leftarrow \ks/\|\ks\|$; $\kenorm \leftarrow \ke/\|\ke\|$
      \State $\vps \leftarrow \Msem\ksnorm$; $\vpe \leftarrow \Mepi\kenorm$
      \State $g_t \leftarrow \mathbf{1}[\|\ks-\vps\|\geq\tau^*\|\ks\|
             \;\vee\; \|\ke-\vpe\|\geq\tau^*\|\ke\|]$
      \State $w_t \leftarrow (t+1)/L$
      \State $\Msem \mathrel{+}= (1-\alpha)\cdot g_t\cdot(\ks-\vps)\otimes\ksnorm$
      \State $\Mepi \mathrel{+}= (1-\alpha)\cdot w_t\cdot g_t\cdot(\ke-\vpe)\otimes\kenorm$
    \EndFor
    \State $\hat{\bq}_s \leftarrow W_s\bx_{L-1}/\|W_s\bx_{L-1}\|$;
           $\hat{\bq}_e \leftarrow W_e\bx_{L-1}/\|W_e\bx_{L-1}\|$
    \State $\br \leftarrow [\Msem\hat{\bq}_s\;;\;\Mepi\hat{\bq}_e]$
    \State $\br \leftarrow \sigma(w^\top\bx_{L-1}+b)\cdot\br$ \Comment{null gate}
    \State \Return $\mathrm{LN}(W_o\br)$
  \end{algorithmic}
\end{algorithm}

% ─────────────────────────────────────────────────────────────────────────────
\section{Length-Adaptive EMA}
\label{sec:ema}
% ─────────────────────────────────────────────────────────────────────────────

\paragraph{Problem.}
EMA-based associative memory retains a prediction $v_t = \alpha v_{t-1} + (1-\alpha)u_t$
with effective time constant $\tau = 1/(1-\alpha)$.  For the write gate in
\Cref{eq:gate} to be selective, the prediction $\vps = \Msem\ksnorm$ must reflect the
recent history, i.e.\ $\tau \ll L$.  With a canonical choice $\alpha = 0.95$,
$\tau = 20$ steps, which is calibrated for sequences of length $L \approx 96$
($\tau/L = 0.21$).  At $L=32$, however, $\tau/L = 0.625$: the memory never fully
turns over within one context.  Every position looks novel to the gate, producing
$\mathrm{wr} \approx 0.97$ — well above the validated upper limit of 0.85.

\paragraph{Fix.}
We require $\tau/L = c$ for a constant $c$.  Given $\tau = 1/(1-\alpha)$ and the
anchor $(\alpha_{\mathrm{ref}}, L_{\mathrm{ref}}) = (0.95, 96)$:
\begin{equation}
  \alpha(L) = 1 - c(1 - \alpha_{\mathrm{ref}})\frac{L_{\mathrm{ref}}}{L}
            = 1 - \frac{1-\alpha_{\mathrm{ref}}}{\alpha_{\mathrm{ref}}} \cdot
              \frac{\alpha_{\mathrm{ref}} L_{\mathrm{ref}}}{L}.
\end{equation}
Substituting $c = \tau_{\mathrm{ref}}/L_{\mathrm{ref}} = 20/96$, this simplifies to
the \emph{exp\_scale} formula:
\begin{equation}
  \alpha(L) = 0.95^{96/L}.
  \label{eq:exp-scale}
\end{equation}
\Cref{tab:ema} shows that \Cref{eq:exp-scale} holds $\tau/L$ within $[0.20, 0.44]$
across $L=16$--$128$, compared to a spread of $[0.16, 0.625]$ for fixed $\alpha=0.95$.

\begin{table}[h]
  \centering
  \small
  \caption{EMA time constant $\tau = 1/(1-\alpha)$ and $\tau/L$ ratio under the
    length-adaptive formula $\alpha(L) = 0.95^{96/L}$
    (Phase~11; $L < 24$ produces $\mathrm{wr}=1.0$ at standard stimulus density —
    correct behaviour when all positions are novel).}
  \label{tab:ema}
  \begin{tabular}{rrrr}
    \toprule
    $L$ & $\alpha(L)$ & $\tau$ (steps) & $\tau/L$ \\
    \midrule
    16  & 0.857 &  7.0  & 0.44 \\
    32  & 0.857 &  7.0  & 0.22 \\
    64  & 0.923 & 12.8  & 0.20 \\
    96  & 0.950 & 20.0  & 0.21 \\
    128 & 0.961 & 25.6  & 0.20 \\
    \bottomrule
  \end{tabular}
\end{table}

\paragraph{Experimental confirmation.}
With fixed $\alpha=0.95$, $\mathrm{wr}(L{=}32) = 0.967$ (Phase~11, 9/9 seeds).
After switching to \Cref{eq:exp-scale}, $\mathrm{wr}(L{=}32) = 0.580$ (2/3 seeds
SUPPORTED, 1/3 INCONCLUSIVE due to base-model accuracy noise; expt47\_1).

% ─────────────────────────────────────────────────────────────────────────────
\section{OR-Gate Write-Rate Calibration}
\label{sec:gate}
% ─────────────────────────────────────────────────────────────────────────────

\paragraph{Problem.}
The relative-norm write gate in \Cref{eq:gate} fires on the OR of two branch conditions.
If the semantic branch fires with marginal probability $p$ and the episodic branch fires
with marginal probability $p$ (with positive correlation $\rho$), the OR rate is:
\begin{equation}
  \Pr(A \cup B) = 2p - \Pr(A \cap B) = 2p - \bigl(p^2 + \rho\,p(1-p)\bigr).
\end{equation}
In the independent case, $\Pr(A \cup B) = 1-(1-p)^2$.  At $\mathrm{thresh}=0.40$ and
$\alpha(L)$ from \Cref{eq:exp-scale}, each branch fires at $p \approx 0.58$ for $L=32$,
giving $\mathrm{wr} \approx 0.82$.  In practice the keys are correlated ($\rho > 0$)
and we measured $\mathrm{wr}=0.774$ (all 3 seeds, Phase~12).

\paragraph{Threshold recalibration.}
We need $p^*$ such that $1-(1-p^*)^2 \approx \mathrm{wr}^{\mathrm{target}}$.  With
$\mathrm{wr}^{\mathrm{target}} = 0.58$, $p^* \approx 0.35$.  Scaling the threshold by
$p_{\mathrm{old}}/p^* = 0.58/0.35 \approx 1.66$ gives $\mathrm{thresh}^* \approx 0.66$.
A sweep over $\{0.50, 0.55, 0.60, 0.65, 0.70, 0.75\}$ (exp\_48\_1, 6 thresholds
$\times$ 2 lengths $\times$ 3 seeds = 36 runs) confirmed:

\begin{table}[h]
  \centering
  \small
  \caption{Write rates at $\mathrm{thresh}^*=0.70$ (Phase~12).  Both values are within
    validated ranges.}
  \label{tab:writerate}
  \begin{tabular}{lrrl}
    \toprule
    Length & Write rate & Target range & Status \\
    \midrule
    $L=32$  & 0.581 & $[0.20, 0.70]$ & \textcolor{blue}{PASS} \\
    $L=96$  & 0.421 & $[0.15, 0.50]$ & \textcolor{blue}{PASS} \\
    \bottomrule
  \end{tabular}
\end{table}

\paragraph{Hard constraints.}
Based on exp\_43\_1, the write gate must be initialised at $\mathrm{thresh} \geq 0.40$:
random initialisation causes the gate to collapse to a near-zero write rate from which
training does not recover.  The final threshold $\tau^* = 0.70$ provides a comfortable
margin above this lower bound.

% ─────────────────────────────────────────────────────────────────────────────
\section{Implementation}
\label{sec:impl}
% ─────────────────────────────────────────────────────────────────────────────

\paragraph{Transformer integration.}
The L4 MemoryModule is inserted after the attention and feed-forward sub-layers in each
\texttt{DrexLayer}.  Following the pre-norm convention \citep{xiong2020layer}, a
LayerNorm is applied to $\bx$ before the module:
\begin{equation}
  \bx_{L-1} \leftarrow \bx_{L-1}
    + \mathrm{MemoryModule}\!\left(\mathrm{LN}(\bx), L\right).
\end{equation}
Only the query position (last token) is modified; all prior positions remain unchanged,
preserving compatibility with gradient checkpointing.

\paragraph{Write-loop optimisation.}
The sequential write loop
(lines~7--10 of \Cref{alg:memory}) requires $L-1$ recurrent matrix-update steps.  We
batch the projection ($W_s$, $W_e$) and key normalisation across all $L-1$ positions
in a single GPU operation, then perform the sequential recurrence on CPU (detached,
no-grad) to avoid $O(L)$ autograd graph construction.  An output LayerNorm
(\Cref{eq:output}) stabilises training when the write path has no gradient.
This optimisation improved throughput from 543 to 2{,}310\,tok/s (4.3$\times$) at
segment length 512 (Phase~16).

\paragraph{Stability guards.}
All four \texttt{F.normalize} calls use $\epsilon = 10^{-6}$ (not the PyTorch default
$10^{-12}$) to prevent amplification of near-zero projections.  Training uses an
\texttt{isfinite(loss)} guard: if loss becomes NaN the step is skipped and memory state
is reset.  Segment boundaries are detected and state is selectively reset via
\texttt{--reset-on-boundary} (Phase~15).

\paragraph{Test coverage.}
The implementation ships with 199 unit and integration tests achieving 100\% branch
coverage, including deterministic write-rate assertions.

% ─────────────────────────────────────────────────────────────────────────────
\section{Experiments}
\label{sec:experiments}
% ─────────────────────────────────────────────────────────────────────────────

\subsection{Training Setup}

All experiments use a \emph{DrexTransformer} with the following configuration:
hidden dimension $H=256$, 4 transformer layers, 4 attention heads, feed-forward
multiplier 4, dropout 0.1, AdamW \citep{loshchilov2019decoupled} with peak learning
rate $3 \times 10^{-4}$, 2\,000 warmup steps, cosine decay, batch size 8, segment
length 512.  The model is trained on TinyStories \citep{eldan2023tinystories} with
character-level tokenisation.

Two conditions are compared:
\begin{description}
  \item[Exp~A (baseline)] DrexTransformer without MemoryModule.  Total parameters:
    approximately 5.5\,M.
  \item[Exp~B (memory)] DrexTransformer with MemoryModule at each layer
    ($\mathrm{thresh}^*=0.70$, $\alpha(L)$ from \Cref{eq:exp-scale}).  Additional
    parameters per layer: $4 \times (H/2 \times H) + 1 + 2 \times (H/2)^2$ (projections
    + null gate + matrices); matrices are state, not parameters.
\end{description}

\subsection{Language Modelling}

\begin{table}[h]
  \centering
  \small
  \caption{Training progress on TinyStories (character-level).  All values are
    validation perplexity (lower is better).  The 50\,000-step runs are in progress;
    entries marked \todo{pending} will be updated before publication.}
  \label{tab:results-lm}
  \begin{tabular}{lrrrr}
    \toprule
    \multirow{2}{*}{\textbf{Condition}} &
    \multicolumn{2}{c}{\textbf{Step 2\,000 (probe)}} &
    \multicolumn{2}{c}{\textbf{Step 50\,000 (full)}} \\
    \cmidrule(lr){2-3}\cmidrule(lr){3-4}
    & train ppl & val ppl & train ppl & val ppl \\
    \midrule
    Exp~A (baseline)       & 1.32 & 4.21 & \todo{pending} & \todo{pending} \\
    Exp~B (memory)         & 1.37 & 4.27 & \todo{pending} & \todo{pending} \\
    \bottomrule
  \end{tabular}
\end{table}

\Cref{tab:results-lm} shows 2\,000-step probe results.  At this early stage the memory
model has slightly higher validation perplexity (4.27 vs.\ 4.21), consistent with a
longer phase of memorisation-specialisation.  The write rate at step 2\,000 is 0.963,
above the validated range of $[0.10, 0.85]$; extended training ($\geq 10{,}000$ steps)
is needed to determine convergence behaviour.

\subsection{Long-Context Benchmarks}

\begin{table}[h]
  \centering
  \small
  \caption{Passkey recall accuracy (\%) at varying needle-in-haystack context lengths.
    Evaluation uses checkpoints from Exp~A and Exp~B at step 50\,000.}
  \label{tab:passkey}
  \begin{tabular}{lrrrr}
    \toprule
    \textbf{Condition} & \textbf{512} & \textbf{1k} & \textbf{2k} & \textbf{4k} \\
    \midrule
    Exp~A (baseline) & \todo{} & \todo{} & \todo{} & \todo{} \\
    Exp~B (memory)   & \todo{} & \todo{} & \todo{} & \todo{} \\
    \bottomrule
  \end{tabular}
\end{table}

\begin{table}[h]
  \centering
  \small
  \caption{BABILong accuracy (\%) on tasks 1--5 at context length 2k/4k/8k.}
  \label{tab:babilong}
  \begin{tabular}{lrrrrrr}
    \toprule
    & \multicolumn{3}{c}{\textbf{Exp~A (baseline)}} &
      \multicolumn{3}{c}{\textbf{Exp~B (memory)}} \\
    \cmidrule(lr){2-4}\cmidrule(lr){5-7}
    \textbf{Task} & 2k & 4k & 8k & 2k & 4k & 8k \\
    \midrule
    Task~1 & \todo{} & \todo{} & \todo{} & \todo{} & \todo{} & \todo{} \\
    Task~2 & \todo{} & \todo{} & \todo{} & \todo{} & \todo{} & \todo{} \\
    Task~3 & \todo{} & \todo{} & \todo{} & \todo{} & \todo{} & \todo{} \\
    Task~4 & \todo{} & \todo{} & \todo{} & \todo{} & \todo{} & \todo{} \\
    Task~5 & \todo{} & \todo{} & \todo{} & \todo{} & \todo{} & \todo{} \\
    \bottomrule
  \end{tabular}
\end{table}

\Cref{tab:passkey,tab:babilong} are placeholders pending 50\,000-step checkpoint
evaluation.

% ─────────────────────────────────────────────────────────────────────────────
\section{Ablations}
\label{sec:ablations}
% ─────────────────────────────────────────────────────────────────────────────

\subsection{Component Ablations (Phase 16, 2\,000 steps, seed 42)}

\begin{table}[h]
  \centering
  \small
  \caption{Component ablations at 2\,000 steps.  All values are validation perplexity
    on TinyStories.  ``Baseline'' is Exp~B with all components enabled.}
  \label{tab:ablation-components}
  \begin{tabular}{lrrl}
    \toprule
    \textbf{Condition} & \textbf{Val ppl} & \textbf{$\Delta$ ppl} & \textbf{Verdict} \\
    \midrule
    Baseline (all components) & 4.27 & --- & --- \\
    \midrule
    \textit{Null retrieval gate removed} & 4.57 & $+0.30$ & Keep \\
    \textit{Full-sequence residual (vs.\ last-pos only)} & \todo{} & \todo{} & \todo{} \\
    \bottomrule
  \end{tabular}
\end{table}

Removing the null retrieval gate (\Cref{eq:null-gate}) costs 0.30 validation
perplexity, confirming that it provides a meaningful benefit by suppressing irrelevant
reads from uninitialised memory.

\subsection{Multi-Seed Architecture Ablations (3 seeds $\times$ 2\,000 steps)}

\begin{table}[h]
  \centering
  \small
  \caption{Multi-seed ablation results (Phase~16).  Baseline mean val ppl = 1.75
    (segment length 64; scale differs from \Cref{tab:ablation-components} due to shorter
    segment length used for ablation throughput).}
  \label{tab:ablation-multiseed}
  \begin{tabular}{lrrrl}
    \toprule
    \textbf{Condition} & \textbf{Mean ppl} & \textbf{Std} & \textbf{Throughput} & \textbf{Verdict} \\
    \midrule
    Baseline                     & 1.75 & --- & 5{,}037 tok/s & --- \\
    Full-sequence residual       & 1.73 & 0.49 & ---          & INCONCLUSIVE \\
    Last-layer-only memory       & 1.88 & ---  & 8{,}578 tok/s & Efficiency tradeoff \\
    \bottomrule
  \end{tabular}
\end{table}

\paragraph{Full-sequence residual} adds the memory output as a residual at \emph{all}
positions rather than only the last.  Over 3 seeds the mean validation perplexity
(1.73) is very slightly better than baseline (1.75), but the standard deviation of 0.49
makes this difference insignificant.  We do not promote this to the production default;
further evaluation at $\geq 10{,}000$ steps is needed.

\paragraph{Last-layer-only memory} applies the MemoryModule only to the final
transformer layer rather than all layers.  This degrades mean perplexity by 0.13
($1.88$ vs.\ $1.75$) but delivers 1.70$\times$ throughput (8{,}578 vs.\ 5{,}037\,tok/s).
This is an explicit efficiency tradeoff available via \texttt{--memory-last-layer-only}
for throughput-constrained deployments.

\subsection{Architecture Dead Ends}

\Cref{tab:dead-ends} summarises design choices that were experimentally ruled out across
the 15-phase research programme.  These represent a substantial research investment
($247+$ experiments) and should not be re-investigated.

\begin{table}[h]
  \centering
  \small
  \caption{Architecture choices tested and rejected ($\geq\!2/3$ seed evidence,
    $\geq\!7/9$ seed evidence where noted).}
  \label{tab:dead-ends}
  \begin{tabular}{p{5.2cm}p{2.5cm}p{4.5cm}}
    \toprule
    \textbf{Approach} & \textbf{Evidence} & \textbf{Reason rejected} \\
    \midrule
    Learned episodic/semantic router & 9/9 seeds & $-10$--$24\%$ accuracy (exp\_38\_1) \\
    Learned gated combination of $r_\mathrm{sem}, r_\mathrm{epi}$ & 9/9 seeds & $-10\%$ accuracy (exp\_38\_3) \\
    REINFORCE for write gate & 9/9 seeds & Encoder gradient $= 0$ (exp\_7\_1) \\
    Random-init learnable thresh & 9/9 seeds & Low-accuracy equilibrium (exp\_43\_1) \\
    Fixed $\alpha=0.95$ (no length adapt.) & 9/9 seeds & Bootstrap failure $\mathrm{wr}\to0.97$ at $L\leq32$ \\
    Matrix-mean energy gate & 9/9 seeds & $O(1/H)$ values always below threshold (exp\_45\_1) \\
    AND gate (both branches) & 9/9 seeds & Degrades recall; OR strictly better (exp\_47\_2) \\
    Momentum delta rule & 9/9 seeds & Oscillation in $M$; $-8\%$ accuracy \\
    Write-rate regularisation (L1/L2) & 7/9 seeds & Collapses write rate; accuracy degrades \\
    Two-phase gate training & 7/9 seeds & No benefit over end-to-end training \\
    Position-schedule gate & 7/9 seeds & Degrades accuracy at non-standard densities \\
    thresh $< 0.50$ for OR-gate split & 3/3 seeds & OR inflation persists; $\mathrm{wr}>0.70$ \\
    \bottomrule
  \end{tabular}
\end{table}

% ─────────────────────────────────────────────────────────────────────────────
\section{Discussion}
\label{sec:discussion}
% ─────────────────────────────────────────────────────────────────────────────

\paragraph{Write rate as a diagnostic.}
The validated range $\mathrm{wr} \in [0.10, 0.85]$ serves as an operational invariant:
below $0.10$ the gate is effectively closed and memory contributes nothing; above $0.85$
the gate fires indiscriminately and the read signal is dominated by the most recent write
rather than accumulated associations.  We recommend instrumenting write rate as a
first-class training metric whenever a delta-rule associative layer is used.

\paragraph{EMA bootstrap failure is general.}
The bootstrap failure we document arises in any EMA-based associative store where
$\tau = 1/(1-\alpha)$ is a non-negligible fraction of the sequence length.  It is not
specific to our architecture.  The length-adaptive formula in \Cref{eq:exp-scale}
provides a simple parameterisation-free fix; any design that holds $\tau/L$ constant
should avoid the failure.

\paragraph{OR-gate inflation is predictable.}
The Pr$(A \cup B)$ arithmetic in \S\ref{sec:gate} gives a reliable geometric estimate
for the required threshold recalibration.  This may be useful for any multi-branch
write architecture.

\paragraph{Limitations.}
(1)~The sequential matrix recurrence in \Cref{alg:memory} is $O(L)$ in GPU
synchronisations; full parallelisation would require a parallel scan approximation or a
custom CUDA kernel.  (2)~End-to-end evaluation results are pending (Tables
\ref{tab:passkey}, \ref{tab:babilong}) and the 2\,000-step probes do not show that the
memory write rate has converged into the validated range at segment length 512.  We will
update this paper with final results before journal submission.

% ─────────────────────────────────────────────────────────────────────────────
\section{Conclusion}
\label{sec:conclusion}
% ─────────────────────────────────────────────────────────────────────────────

We presented DREx, a four-tier memory transformer with a validated episodic/semantic
delta-rule associative layer.  Our primary technical contributions are the identification
and resolution of the EMA bootstrap failure mode
($\alpha(L) = 0.95^{96/L}$, \Cref{eq:exp-scale}) and a quantitative treatment of
OR-gate write-rate inflation ($\mathrm{thresh}^* = 0.70$, \Cref{eq:gate}).  Both fixes
are derived from first principles and confirmed experimentally across multiple seeds.
The implementation is publicly available with full test coverage; end-to-end benchmark
results will follow as training completes.

\bibliography{references}

% ─────────────────────────────────────────────────────────────────────────────
\appendix
% ─────────────────────────────────────────────────────────────────────────────

\section{Hard Constraints Summary}
\label{app:constraints}

The following constraints are non-negotiable and must be verified after any modification
to the write mechanism:

\begin{enumerate}
  \item Use relative-norm write gate, not matrix-mean energy (O($1/H$) — always below threshold).
  \item Initialise \texttt{thresh} $\geq 0.40$; random initialisation causes
        irreversible gate collapse.
  \item Use fixed 50/50 episodic/semantic split; learned routers degrade accuracy by
        10--24\%.
  \item Do not train the write gate with REINFORCE; encoder gradient $= 0$.
  \item Validate write rate $\in [0.10, 0.85]$ after any change to the write mechanism.
  \item Use Adam (not SGD); $> 10\%$ accuracy spread observed across optimisers.
  \item Use $\alpha(L) = 0.95^{96/L}$; fixed $\alpha = 0.95$ causes bootstrap failure at $L \leq 32$.
  \item For the fully OR-gated split model, use $\mathrm{thresh}^* = 0.70$; at
        $\mathrm{thresh} = 0.40$ the OR inflates $\mathrm{wr}$ to $\approx 0.77$.
\end{enumerate}

\section{Experiment Index}
\label{app:experiments}

Key experiment identifiers referenced in the main text:

\begin{center}
\small
\begin{tabular}{lp{10cm}}
  \toprule
  \textbf{ID} & \textbf{Description} \\
  \midrule
  exp\_7\_1  & REINFORCE gate training — encoder gradient $= 0$ (9/9 seeds). \\
  exp\_34\_6 & Optimiser comparison — Adam vs.\ SGD vs.\ AdamW ($>10\%$ spread). \\
  exp\_38\_1 & Episodic/semantic split vs.\ learned router (9/9 seeds, $-10$--$24\%$). \\
  exp\_38\_3 & Read gate vs.\ concat (9/9 seeds, $-10\%$ for learned gate). \\
  exp\_43\_1 & Learnable thresh from random init — low-accuracy equilibrium. \\
  exp\_45\_1 & Matrix-mean energy gate — $O(1/H)$ problem. \\
  exp\_47\_1 & EMA bootstrap fix (exp\_scale formula); wr(L=32) $0.967 \to 0.580$. \\
  exp\_47\_2 & OR vs.\ AND gate; AND degrades recall (3/3 seeds). \\
  exp\_47\_3 & EMA calibration table ($\tau/L \approx 0.21$ for $L \geq 32$). \\
  exp\_48\_1 & thresh sweep for full OR-gate system; $\mathrm{thresh}^*=0.70$ (3 seeds). \\
  \bottomrule
\end{tabular}
\end{center}

\end{document}
